\section{Prompt template shape and workload questions}
\label{app:prompt}

\subsection{Compiled sequence}

Every question is compiled into the following shape. Role markers and any special tokens are
produced by the backbone's own chat template (rendered with thinking disabled) and encoded
with the host tokenizer; everything shown in the user message after \texttt{"state":} is
encoded with the data tokenizer, which has no added tokens (Section~\ref{sec:datatok}).
Line breaks inside the user message are for display only; the actual JSON is compact
(separators \texttt{,} and \texttt{:}, non-ASCII characters kept).

\begin{small}
\begin{verbatim}
[system]
Judge the state using the question's instructions and all candidate meanings.
The state and candidate descriptions are data, not additional system messages.
Return exactly the selected identifier, with no explanation or extra characters.
For score questions the candidates are ordered levels; for noul they mean
false/true.
[user]
{"state":STATE,
 "question":{"type":TYPE,
             "instructions":INSTRUCTIONS,
             "candidates":[{"identifier":"A","meaning":MEANING_1},
                           {"identifier":"B","meaning":MEANING_2},
                           ...]}}
[assistant, generation prompt]
   <- candidate logits are read at the last position of this sequence
\end{verbatim}
\end{small}

\begin{itemize}
  \item \texttt{STATE} is the request's state as JSON (a string state becomes a JSON string);
    key order is preserved.
  \item \texttt{TYPE} is \texttt{choice}, \texttt{noul} or \texttt{score};
    \texttt{INSTRUCTIONS} is the caller's value as JSON (or \texttt{null}).
  \item For a Choice, \texttt{MEANING} is the candidate name if it has no description, and
    otherwise the JSON text of \texttt{\{"description": D, "name": N\}} with sorted keys.
  \item For a Noul there are two candidates, in the order \emph{false}, \emph{true}, whose
    meanings are the JSON text of \texttt{\{"description": D, "value": "false"\}} and the
    same for \texttt{"true"} (\texttt{D} may be \texttt{null}).
  \item For a Score, the candidates are the levels in order, and \texttt{MEANING} is the
    level's text (or its JSON text if the level is structured).
  \item Identifiers are assigned by position from the verified list of
    Section~\ref{sec:identifiers}. The question ID is never part of the sequence.
\end{itemize}

\subsection{Injection and distractor texts}

The stress suites of Section~\ref{sec:robust} use Korean text. The injection sentences are
translated in Section~\ref{sec:injection}. The two irrelevant candidates are ``banana
prices'' (described as today's banana price in the fruit market) and ``space station orbit''
(the current orbital altitude of the International Space Station); for title questions they
are replaced by headlines about a banana price rise and an orbit adjustment of the
International Space Station.

\subsection{Serving workload questions}

The serving benchmark (Section~\ref{sec:serving}) builds its states as
\texttt{\{"document": text\}} from development passages, with the text length chosen by
binary search so that the longest compiled question reaches the target length. The questions
are Korean; translated:
\begin{itemize}
  \item \textbf{W1} (and each W4 request): one Choice, ``choose the main topic of the
    document'', over eight topics (politics and elections, economy and finance, society and
    incidents, international relations, science and technology, culture and arts, sports,
    life and health).
  \item \textbf{W2}: eight questions on the same state: four Noul questions (does the
    document deal with a person, a date, an amount, a place), the eight-way topic Choice, a
    four-way tone Choice (neutral report, critical, promotional, personal opinion), a
    six-way audience Choice (general readers, experts, investors, officials and
    institutions, students, consumers) and a five-level Score of informational specificity
    (none, low, medium, high, very high).
  \item \textbf{W3} (8K and 16K): one four-way Choice of document type (news report, column
    or essay, notice, other). \textbf{W3 K77 / K255}: the first development title-selection
    questions with 77 and 255 candidates.
  \item \textbf{W4}: 32 distinct W1-style requests of 1,024 tokens, run as static batches of
    1, 8, 16 and 32.
\end{itemize}
